\scanpage{190}

à l'ens des parties d'un ens. fini, est spécial \unsure{aux} croisements normaux ce n'est déjà plus le cas pour
\figure{\begin{tikzpicture}[baseline=-.3ex,scale=.9]
\draw (0.02,.09) -- (.18,.69);
\draw (-.10,-.05) to[out=200,in=15] (-.56,-.21);
\draw (.10,-.06) -- (.48,-.25);
\draw (0,0) circle[radius=.035];
\end{tikzpicture}}, où le nb des strates contenant $\bullet$ est (en comptant $\bullet$ lui-$=$) de $7$, qui n'est pas une puissance de deux...

\uline{NB} Pour l'instant, on n'a pas fait usage de l'hypothèse \inserted{\unsure{du moins|des voisinages}} de locale trivialité de p.\ 86 e). Elle servira plus loin, pour expliciter en termes de fibrations loc.\ triviales certains voisinages tubulaires, \deleted{qui \unsure{seront à l'origine}} \inserted{jouant le rôle} des ``topos élémentaires'' de la section précédente.

\margin{189}
